%==========================================================
\section{Diagrama de afinidade}\label{sec:afinidade}

O diagrama de afinidade organiza as notas de interpretação das duas sessões (A0-01 a A0-15 e U0-01 a U0-21) por semelhança de significado, de baixo para cima, para revelar os temas comuns a toda a população~\cite[cap.~9]{beyer1998}. As notas agrupam-se em grupos de 4--6, cada grupo recebe um rótulo azul escrito na voz do utilizador, os grupos azuis juntam-se sob rótulos cor-de-rosa e estes sob rótulos verdes.

\pendente{a definir}{%
\begin{itemize}
  \item Construir o diagrama a partir das 36 notas de interpretação (Secção~\ref{sec:recolha} e Apêndice~\ref{anx:notas}), sem incluir as notas do tipo OP sem caso concreto.
  \item Incluir a figura do diagrama (ou uma tabela hierárquica verde $\rightarrow$ rosa $\rightarrow$ azul $\rightarrow$ notas).
  \item Indicar para cada rótulo azul as notas que o sustentam, para manter a rastreabilidade.
  \item Se for usada IA para propor agrupamentos, guardar a versão da IA e a versão final da equipa para o Apêndice~\ref{anx:ia}.
\end{itemize}}

% Exemplo de estrutura em tabela, a substituir pelo diagrama real:
% \begin{table}[!htbp]\centering\small
% \begin{tabularx}{\textwidth}{@{}p{3.2cm}p{4cm}L@{}}
% \toprule
% \textbf{Verde} & \textbf{Azul (voz do utilizador)} & \textbf{Notas} \\ \midrule
% Custo real da viagem & «Só sei quanto custa quando já estou a pagar» & U0-07, U0-08 \\
% \bottomrule
% \end{tabularx}
% \caption{Diagrama de afinidade.}\label{tab:afinidade}
% \end{table}
